\section{Discussion and Future Work}
\label{sec:discussion}

\paragraph{Writing \(\Phi\).}
Decoupling search from \(P_1\) moves the specification burden onto the
developer. That cost has limited adoption of contract-based tools.
A contract inferred from the baseline tends to freeze incidental choices,
such as a capacity of exactly 1000, and removes the slack the problem
depends on. Contracts have to come from intent, not from the code.
Translating a goal-oriented requirements model into \(\mathit{Pre}\),
\(\mathit{Post}\), and \(\mathcal{I}\) is future work. One route is a
KAOS refinement of a goal tree down to atomic requirements in a formal
specification language~\cite{vanlamsweerde2001gore}.

\paragraph{Under-Specification.}
The legal set is the slack in \(\Phi_{\text{func}}\), including slack the
developer did not intend. A FIFO ring whose contract omits order admits a
faster queue that reorders packets, and \(V\) will prove it. Checking a
proposed contract against the baseline catches contracts that are too
strong, not contracts that are too weak~\cite{misu2026specharness}. A missing
requirement is a proved program the developer did not want. When a
deployment or a rejected candidate surfaces a gap, that requirement
should become a contract clause that later compilations re-check, so the
same regression cannot return unless the requirements change. Completing
\(\Phi\) that way is the developer's cycle around the compiler, not a
pass of Section~\ref{sec:compiler}. The same holds for a cost that is
missing from \(\mathcal{B}\) or \(\mathcal{O}\): concurrent false sharing
on a padded ring cannot reject a candidate until it is named. Closing
gaps permanently is the central open problem for intent compilation.

\paragraph{Temporarily Relaxing Constraints.}
Loosening \(\mathit{Pre}\) or \(\mathit{Post}\) enlarges
\(\mathrm{Sat}(\Phi_{\text{func}})\), so the search can admit a program the
tight contract would reject. Putting those clauses back makes that program
illegal, and \(V\) has to run again. Loosening \(\mathcal{B}\) does not
change legality: a candidate admitted under a looser memory ceiling can
still satisfy \(\Phi_{\text{func}}\) after the tight ceiling is restored,
and no observable test fails. The compiler also bounds how much of
\(\mathrm{Gen}\) it explores and how long \(V\) may run; loosening those
changes which candidates finish in a compile, not
\(\mathrm{Sat}(\Phi_{\text{func}})\). Restoring a temporary relaxation is
future work: re-running \(V\), re-checking \(\mathcal{B}\), or re-searching
under the tight bounds.

\paragraph{Unbounded Verification.}
\textsc{IntC}'s claim is \(\models_K\). An unbounded proof of a stateful
component is the invariant-and-abstraction obligation Ivy discharges by
restructuring into a decidable
fragment~\cite{padon2016ivy,padon2017paxosepr}. Carrying that obligation
on a candidate, rather than stopping at \(K\), is future work.

\paragraph{Where to Spend Compilation Time.}
\textsc{IntC} runs \(V\) before it measures \(\mathcal{O}\), so a
candidate can consume a timeout before a cheap profile would have
dropped it. Other schedules over the same \(\mathrm{Gen}\) include
estimating \(\mathcal{O}\) first, running a cheap \(\models_K\) before a
longer query, and keeping candidates and proofs as artifacts so a later
compile reuses them.

\paragraph{A Score that Predicts \(V\).}
\textsc{IntC}'s TheoryLower score orders descent, and against CBMC time it
has Spearman's \(\rho = 0.11\). A score that prefers the rewrite expected
to shorten the query has to be fit to this \(V\), and predicting SMT and
verifier runtime from goal or program structure is already
studied~\cite{healy2017where4,leeson2022graves}.

\paragraph{Language Models as Compiler Inputs.}
\label{sec:discussion_lm}
A model can be the untrusted \(\mathrm{Gen}\) of
Theorem~\ref{thm:restricted} at compile
time~\cite{taneja2025llmvectorizer,fang2026llmveriopt,jiang2026segabench},
or, for a deterministic compile, help write the catalogs, TheoryLower
scores, and benchmarks an instance enumerates. It can also help with the
specification burden of writing \(\Phi\) by drafting
\(\Phi_{\text{func}}\) from a prose
requirement~\cite{ma2025specgen,ugare2026fun2spec} as a candidate for
\(V\), and by helping write the baseline \(P_1\). \textsc{IntC} is the prototype
for a C instance; a model can help instantiate those inputs for another
language.

\paragraph{Language Models as Developers.}
Professional developers commonly use AI tools on coding tasks: 51\%
report daily use in the 2025 Stack Overflow
survey~\cite{stackoverflow2025ai}, including multi-turn conversations on
real GitHub work~\cite{hao2024chatgpt}. Later rewrites in that setting
often break requirements that earlier turns had already
met~\cite{huang2026regression,yan2026slump}.
\textsc{IntentCompile} takes those requirements as a checkable
\(\Phi\) and the program as \(P_1\), which a model can
fill~\cite{lahiri2026intent}. The next rewrite is then a compile against
\(\Phi\). Named requirements change when intent changes, which is less
often than the program is rewritten, so a human can maintain \(\Phi\)
while the model rewrites \(P\).

\section{Conclusion}
\label{sec:conclusion}

An as-if compiler preserves the source, including incidental choices the
requirements do not fix, so it may omit a cheaper legal program.
\textsc{IntentCompile} takes a baseline and a checkable contract,
generates candidates, and admits those a sound verifier proves. It ranks
the admitted set by \(\mathcal{O}\) under the same lowering backend that
compiled \(P_1\), and emits \(P_1\) when no admitted candidate is cheaper.
\textsc{IntC} is a prototype C instance of that restricted compiler.

On the motivating example, \textsc{IntC} admits the 1024-cell ring that
MSVC~\texttt{/O2} will not emit from a 1000-cell remainder. On extracts
whose contracts name a benchmark, the emit is cheaper on \(\mathcal{O}\)
than \(P_1\) or is \(P_1\). The catalogs, the bound \(K\), and writing
\(\Phi\) from intent are the limits of \textsc{IntC}.
